% Copyright (c) 2026 Benjamin Stanley Frohman.
% Paper text: CC BY 4.0. Quoted Lean code: Apache-2.0 (BenFrohman/frohmanian-tether-sdiff).
\documentclass[11pt]{article}
\usepackage[margin=1in]{geometry}
\usepackage{fontspec}
\setmonofont{DejaVu Sans Mono}[Scale=0.78]
\newfontfamily\bigsqcupfont{Noto Sans Mono}[Scale=0.78]
\XeTeXinterchartokenstate=1
\newXeTeXintercharclass\bigsqcupclass
\XeTeXcharclass"2A06=\bigsqcupclass
\XeTeXinterchartoks 0 \bigsqcupclass={\begingroup\bigsqcupfont}
\XeTeXinterchartoks 4095 \bigsqcupclass={\begingroup\bigsqcupfont}
\XeTeXinterchartoks \bigsqcupclass 0={\endgroup}
\XeTeXinterchartoks \bigsqcupclass 4095={\endgroup}
\usepackage{amsmath,amssymb,amsthm}
\usepackage{fvextra}
\usepackage{xcolor}
\usepackage[hidelinks]{hyperref}
\usepackage{enumitem}
\fvset{fontsize=\small,frame=single,framesep=2mm,breaklines=true,breakanywhere=true,fontshape=n}
\sloppy\emergencystretch=4em
\usepackage{seqsplit}
\newcommand{\R}{\mathbb{R}}
\newcommand{\lean}[1]{\texttt{#1}}
\newcommand{\commit}{\texttt{\seqsplit{1ea649e495b5e06fafec75a05bb8ec07f13dae80}}}
\newcommand{\prcommit}{\texttt{\seqsplit{ad546216d56b5ab7514f1ed20fe168744ff9cf74}}}
\newcommand{\hash}[1]{\texttt{\seqsplit{#1}}}
\newtheorem{theorem}{Theorem}
\theoremstyle{definition}
\newtheorem{definition}{Definition}
\theoremstyle{remark}
\newtheorem*{remark}{Remark}

\title{Local Kato Theory for the Incompressible Navier--Stokes Equations on $\R^3$, Kernel-Checked in Lean~4:\\
A Maximal Solution with an $H^3$ Blow-up Alternative for $\nu>0$\\
\large (Not Global Regularity; Not Fefferman's Problem~(A))}
\author{Benjamin Stanley Frohman\\
\small Independent Researcher, Texas, United States\\
\small ORCID \href{https://orcid.org/0009-0006-7068-3718}{0009-0006-7068-3718}}
\date{17 November 2026 (America/Chicago) --- preprint, not peer-reviewed\\ \small DOI \href{https://doi.org/10.5281/zenodo.24302258}{10.5281/zenodo.24302258}}

\begin{document}
\maketitle

\begin{abstract}
For every viscosity $\nu>0$ and every smooth, divergence-free initial field $u_0$ on $\R^3$ satisfying Fefferman's decay condition~(4), the Lean~4 proposition \lean{KatoMaximalSolution ν u₀} is proved and checked by the Lean kernel: there exist a time set $S$, a velocity $u$, a pressure $p$ and a constant $C_k\ge 0$ such that $(u,p)$ is a smooth Navier--Stokes solution from $u_0$ (in the class \lean{IsFeffermanSolutionOn}) on every $[0,T)\subseteq S$; $\|u(t)\|_{H^3}^2$ is finite on $S$ and continuous on $S$; at every $s\in S$ the map $t\mapsto\|u(t)\|_{H^3}^2$ has a right derivative $d\le C_k\,\|\nabla u(s)\|_\infty\,\|u(s)\|_{H^3}^2$; and either $S=[0,\infty)$, or $S=[0,T^*)$ with $T^*>0$ and $\|u(t)\|_{H^3}^2$ unbounded on $S$. This is local Kato theory. It is not a global regularity theorem, it does not prove Fefferman's Problem~(A), and it does not solve any open problem. Two further kernel-checked statements are conditional: \lean{R3BKMContinuation ν} keeps a uniform vorticity ceiling as a hypothesis, and \lean{clay\_A\_fefferman\_of\_rate} keeps \lean{FrohmanianTetherRateSmooth} as a hypothesis, which is not proved. Every statement is quoted verbatim from \lean{Scratch/KatoMaximal.lean} and \lean{Scratch/KatoContinuation.lean} at commit \commit{}, whose Lean sources are identical to \prcommit{} (the merge of pull request~\#12), and was built by \lean{lake build} and \lean{lake build Scratch} on GitHub Actions run 44047249311 with the repository's own \lean{lake-manifest.json} (sha256 \lean{7d1ab35f…0bc0955}) and elan 4.2.4; \lean{\#print axioms} reports only \lean{propext}, \lean{Classical.choice} and \lean{Quot.sound} for each.
\end{abstract}

\noindent\textbf{Copyright (c) 2026 Benjamin Stanley Frohman.} Paper: CC BY 4.0. Lean code: Apache-2.0.

\section{Scope}
\label{sec:scope}
This paper reports exactly the Lean statements listed in Sections~\ref{sec:main}--\ref{sec:cond}, as checked by the Lean kernel in the build described in Section~\ref{sec:repro}. Nothing else is claimed.
\begin{itemize}[leftmargin=*]
\item \textbf{Proved (unconditional, for $\nu>0$ and Fefferman data):} Theorem~\ref{thm:T1} (\lean{KatoMaximalSolution}) and Theorems~\ref{thm:T2}--\ref{thm:T9} (properties of the realized maximal solution on its maximal time set).
\item \textbf{Proved, conditional in content:} Theorem~\ref{thm:T10}, whose conclusion \lean{R3BKMContinuation ν} is itself an implication with a uniform vorticity ceiling as its premise; Theorem~\ref{thm:T11}, whose hypothesis \lean{FrohmanianTetherRateSmooth} is not proved.
\item \textbf{Not proved:} global regularity for the Navier--Stokes equations on $\R^3$; Fefferman's Problem~(A) (\lean{Clay\_A\_Fefferman}); a uniform bound on $\|\omega\|_\infty$; \lean{FrohmanianTetherRateSmooth}. No Millennium Prize problem is solved.
\end{itemize}

\section{Setting and notation}
\label{sec:setting}
Space is $\R^3$, in Lean \lean{T3 := EuclideanSpace ℝ (Fin 3)}; a velocity field is \lean{VelocityField := T3 → EuclideanSpace ℝ (Fin 3)}, a pressure field is \lean{T3 → ℝ}, and time-dependent fields are functions \lean{ℝ → VelocityField}, \lean{ℝ → PressureField}. The operators are defined in \lean{FrohmanianTether/NS\_Equations.lean}: \lean{convective u v x = (fderiv ℝ v x) (u x)}, written $(u\cdot\nabla)v$; \lean{laplacian}, written $\Delta$, is the trace of the Hessian; \lean{div}, written $\operatorname{div}$, is $\sum_i\partial_i u_i$; \lean{pressureGradient p}, written $\nabla p$, is the Hilbert-space gradient. ``Smooth'' is \lean{ContDiff ℝ ∞}. The definitions below are quoted verbatim from the source at commit \commit.

\begin{definition}[Fefferman decay (4)] \lean{Scratch/RepairedTetherRate.lean:33--36}: $|\partial_x^\alpha u_0(x)|\le C_{\alpha K}(1+|x|)^{-K}$ for every multi-index $\alpha$ and every real $K$.
\VerbatimInput{lean/DFD.lean}
\end{definition}

\begin{definition}[Local Fefferman-class solution on $\R^3\times[0,T)$] \lean{Scratch/RepairedTetherRate.lean:66--75}: $u$ and $p$ are jointly $C^\infty$ on $\R^3\times[0,T)$; for $t\in[0,T)$ and $x\in\R^3$, $\partial_t^+u+(u\cdot\nabla)u+\nabla p=\nu\Delta u$ (one-sided time derivative within $[0,\infty)$) and $\operatorname{div}u=0$; $u(\cdot,0)=u_0$; and $\int_{\R^3}|u(x,t)|^2\,dx\le\int_{\R^3}|u_0(x)|^2\,dx$ for $t\in[0,T)$, as lower Lebesgue integrals.
\VerbatimInput{lean/DIFS.lean}
\end{definition}

\begin{definition}[$H^3$ energy and $\|\nabla u\|_\infty$] \lean{Scratch/R3BKMContinuation.lean:44--55}. Throughout this paper $\|f\|_{H^m}^2$ denotes \lean{hsSq m f}, the real value of
$\sum_{k\le m}\sum_{\alpha:\{0,\dots,k-1\}\to\{1,2,3\}}\int_{\R^3}|\partial_{\alpha_0}\cdots\partial_{\alpha_{k-1}}f(x)|^2\,dx$ (ordered tuples of directions), and $\|\nabla f\|_\infty$ denotes \lean{gradSup f}$\,=\sup_x\|Df(x)\|$ as a real supremum.
\VerbatimInput{lean/DsobE.lean}
\VerbatimInput{lean/DhsSq.lean}
\VerbatimInput{lean/DgradSup.lean}
\end{definition}

\begin{remark}
In Mathlib a real supremum of a family that is not bounded above equals $0$, and \lean{ENNReal.toReal ⊤ = 0}. Hence \lean{gradSup f} and \lean{hsSq m f} are the textbook quantities only when $\nabla f$ is bounded and $\|f\|_{H^m}^2<\infty$ respectively. Theorem~\ref{thm:T1} (Definition~\ref{def:kato}(ii)) asserts $\|u(t)\|_{H^3}^2<\infty$ on $S$. None of the theorems below asserts that $\nabla u(s)$ is bounded.
\end{remark}

\begin{definition}[Kato maximal solution]\label{def:kato} \lean{Scratch/R3BKMContinuation.lean:88--96}. \lean{KatoMaximalSolution ν u₀} states that there exist a time set $S$, fields $(u,p)$ and a constant $C_k\ge 0$ such that:
\begin{enumerate}[label=(\roman*)]
\item for every $T$ with $[0,T)\subseteq S$, $(u,p)$ is a local Fefferman-class solution from $u_0$ on $[0,T)$;
\item $\|u(t)\|_{H^3}^2<\infty$ for every $t\in S$; the map $t\mapsto\|u(t)\|_{H^3}^2$ is continuous on $S$; and at every $s\in S$ it has a right derivative $d$ (derivative within $[s,\infty)$) with $d\le C_k\,\|\nabla u(s)\|_\infty\,\|u(s)\|_{H^3}^2$;
\item either $S=[0,\infty)$, or $S=[0,T^*)$ for some $T^*>0$ and for every $B\in\R$ there is $t\in S$ with $\|u(t)\|_{H^3}^2>B$.
\end{enumerate}
\VerbatimInput{lean/DKMS.lean}
\end{definition}

\section{The main theorem}
\label{sec:main}
All theorems in Sections~\ref{sec:main} and~\ref{sec:smax} live in \lean{section Lift} of \lean{Scratch/KatoMaximal.lean}, whose variables are
\VerbatimInput{lean/Dvar.lean}
so each theorem takes the hypotheses $u_0\in C^\infty$ (\lean{hf}), Fefferman decay (\lean{hdec}) and $\nu>0$ (\lean{hν}), together with the binders shown. Theorem~\ref{thm:T1} is preceded in the source by \lean{include hf hdec hν in}; it is applied as \lean{katoMaximalSolution\_of\_feffermanDecay hs hdec hν hd} in \lean{Scratch/KatoContinuation.lean:44}.

\begin{theorem}[Kato maximal solution for Fefferman data]
\label{thm:T1}
Let $\nu>0$, and let $u_0$ be smooth, divergence-free, and satisfy Fefferman decay~(4). Then \lean{KatoMaximalSolution}\ $\nu$ $u_0$ holds in the sense of Definition~\ref{def:kato}: there exist a time set $S$, fields $(u,p)$ and a constant $C_k\ge 0$ satisfying (i)--(iii) of that definition. This is a maximal Fefferman-class solution with an $H^3$ blow-up alternative; it does not assert global regularity.

\noindent\lean{KatoMaximal.katoMaximalSolution\_of\_feffermanDecay}, \lean{Scratch/KatoMaximal.lean:2103--2104}:
\VerbatimInput{lean/T1.lean}
\end{theorem}

The witness in the proof is $S=$\,\lean{Smax hf hdec hν (by norm\_num : 2 ≤ 2)}, $u=$\,\lean{physicalMaximal hf hdec hν}, $p=$\,\lean{presMaximal hf hdec hν}. Theorems~\ref{thm:T2}--\ref{thm:T9} are stated for these objects. \lean{Smax} is defined at \lean{Scratch/KatoMaximal.lean:561--565}:
\VerbatimInput{lean/DinMax.lean}
\VerbatimInput{lean/DSmax.lean}

\section{Properties on the maximal time set}
\label{sec:smax}
Write $S_{\max}^{(k)}$ for \lean{Smax hf hdec hν hk}, $u(t)$ for \lean{physicalMaximal hf hdec hν t} and $p(t)$ for \lean{presMaximal hf hdec hν t}.

\begin{theorem}[$H^3$ Kato--Ponce bound on $S_{\max}$]
\label{thm:T2}
If $\operatorname{div}u_0=0$, there is $C_k\ge 0$ such that at every $s\in S_{\max}^{(2)}$ the map $t\mapsto\|u(t)\|_{H^3}^2$ has a right derivative $d$ at $s$ with $d\le C_k\,\|\nabla u(s)\|_\infty\,\|u(s)\|_{H^3}^2$.
\lean{Scratch/KatoMaximal.lean:2071--2078}:
\VerbatimInput{lean/T2.lean}
\end{theorem}

\begin{theorem}[Vanishing-flux package at each time]
\label{thm:T3}
If $\operatorname{div}u_0=0$ and $t\in S_{\max}^{(2)}$, then $(u(t),p(t))$ satisfies \lean{H3IBPVanishingFlux}: for every order $k\le 3$ and every $\alpha$, the integrability and zero-flux conditions of \lean{WordIBP} (quoted below) hold.
\lean{Scratch/KatoMaximal.lean:2036--2039}:
\VerbatimInput{lean/T3.lean}
\VerbatimInput{lean/DH3.lean}
\VerbatimInput{lean/DWord.lean}
\end{theorem}

\begin{theorem}[Fefferman-class solution on every $[0,T)\subseteq S_{\max}$]
\label{thm:T4}
If $\operatorname{div}u_0=0$ and $[0,T)\subseteq S_{\max}^{(2)}$, then $(u,p)$ is a local Fefferman-class solution from $u_0$ on $[0,T)$.
\lean{Scratch/KatoMaximal.lean:1809--1812}:
\VerbatimInput{lean/T4.lean}
\end{theorem}

\begin{theorem}[Momentum equation]
\label{thm:T5}
If $\operatorname{div}u_0=0$, $[0,T)\subseteq S_{\max}^{(2)}$, $t\in[0,T)$ and $x\in\R^3$, then $\partial_t^+u(x,t)+(u(t)\cdot\nabla)u(t)(x)+\nabla p(t)(x)=\nu\,\Delta u(t)(x)$, with the time derivative taken within $[0,\infty)$.
\lean{Scratch/KatoMaximal.lean:1626--1633}:
\VerbatimInput{lean/T5.lean}
\end{theorem}

\begin{theorem}[Energy inequality]
\label{thm:T6}
If $\operatorname{div}u_0=0$ and $[0,T)\subseteq S_{\max}^{(2)}$, then $\int_{\R^3}|u(x,t)|^2\,dx\le\int_{\R^3}|u_0(x)|^2\,dx$ for every $t\in[0,T)$ (lower Lebesgue integrals).
\lean{Scratch/KatoMaximal.lean:1793--1796}:
\VerbatimInput{lean/T6.lean}
\end{theorem}

\begin{theorem}[Continuity of the $H^3$ energy]
\label{thm:T7}
For every $k\ge 2$, $t\mapsto\|u(t)\|_{H^3}^2$ is continuous on $S_{\max}^{(k)}$.
\lean{Scratch/KatoMaximal.lean:1193--1195}:
\VerbatimInput{lean/T7.lean}
\end{theorem}

\begin{theorem}[$H^3$ blow-up alternative]
\label{thm:T8}
For every $k\ge 2$, either $S_{\max}^{(k)}=[0,\infty)$, or there is $T^*>0$ with $S_{\max}^{(k)}=[0,T^*)$ and, for every $B\in\R$, some $t\in[0,T^*)$ with $\|u(t)\|_{H^3}^2>B$.
\lean{Scratch/KatoMaximal.lean:1229--1234}:
\VerbatimInput{lean/T8.lean}
\end{theorem}

\begin{theorem}[Shape of the maximal time set]
\label{thm:T9}
For every $k\ge 2$, either $S_{\max}^{(k)}=[0,\infty)$ or $S_{\max}^{(k)}=[0,T^*)$ for some $T^*>0$.
\lean{Scratch/KatoMaximal.lean:1052--1054}:
\VerbatimInput{lean/T9.lean}
\end{theorem}

\section{Two conditional statements}
\label{sec:cond}
Both theorems below are in \lean{Scratch/KatoContinuation.lean}. They use the following definitions from \lean{Scratch/RepairedTetherRate.lean}, quoted verbatim.
\VerbatimInput{lean/DMw.lean}
\VerbatimInput{lean/DContAt.lean}
\VerbatimInput{lean/DCont.lean}
\VerbatimInput{lean/DFGS.lean}
\VerbatimInput{lean/DClay.lean}
\VerbatimInput{lean/DRate.lean}

\begin{theorem}[Classical BKM continuation on $\R^3$]
\label{thm:T10}
For every $\nu>0$, \lean{R3BKMContinuation}\ $\nu$ holds: for every smooth, divergence-free $u_0$ satisfying~(4), if one ceiling $Y$ bounds $\|\omega\|_\infty$ along every local Fefferman-class solution from $u_0$, that is, $\|\omega(t)\|_\infty$ is finite and $\|\omega(t)\|_\infty\le Y$ for every local Fefferman-class solution $(u,p)$ from $u_0$ on $[0,T)$ and every $t\in[0,T)$, then $u_0$ has a global Fefferman solution: smooth functions $p$ and $u$ on $\R^3\times[0,\infty)$ that satisfy (1), (2), (3), (6) and (7) of Fefferman's problem description. A uniform vorticity ceiling remains a hypothesis; this is not a regularity theorem.

\noindent\lean{KatoContinuation.r3BKMContinuation\_of\_pos}, \lean{Scratch/KatoContinuation.lean:41}:
\VerbatimInput{lean/T10.lean}
\end{theorem}

\begin{theorem}[Conditional reduction to the Smooth rate alone]
\label{thm:T11}
Assume that for every $\nu>0$ and every Fefferman datum $u_0$ (smooth, divergence-free, satisfying~(4)), \lean{FrohmanianTetherRateSmooth}\ $u_0$ $\nu$ holds. Then \lean{Clay\_A\_Fefferman} holds. The rate hypothesis is not proved in this repository, so this does not establish \lean{Clay\_A\_Fefferman}.

\noindent\lean{KatoContinuation.clay\_A\_fefferman\_of\_rate}, \lean{Scratch/KatoContinuation.lean:49--52}:
\VerbatimInput{lean/T11.lean}
\end{theorem}

\section{What is not proved}
\label{sec:notproved}
\begin{itemize}[leftmargin=*]
\item Global existence of smooth solutions on $\R^3\times[0,\infty)$ for general Fefferman data. Theorem~\ref{thm:T1} (Definition~\ref{def:kato}(iii)) leaves both alternatives open: no statement here excludes $S=[0,T^*)$.
\item \lean{Clay\_A\_Fefferman}, i.e.\ Fefferman's Problem~(A). Theorem~\ref{thm:T11} derives it only from \lean{FrohmanianTetherRateSmooth} for all Fefferman data, which is not proved.
\item A uniform ceiling on $\|\omega\|_\infty$ along local Fefferman-class solutions. It is the premise inside \lean{R3BKMContinuationAt}, and Theorem~\ref{thm:T10} does not discharge it.
\item Boundedness of $\nabla u(s)$ on $S_{\max}$ (see the Remark in Section~\ref{sec:setting}).
\end{itemize}

\section{Truth Gate status}
\label{sec:gate}
Theorems~\ref{thm:T1}, \ref{thm:T10} and~\ref{thm:T11} are covered by the author-approved Truth Gate locks \lean{kato}, \lean{bkm}, \lean{clayA} in \lean{Scratch/FeffermanSemantics.lean} (lock record \lean{fefferman-clay-a-conditional} in \lean{TRUTH\_GATE.lock}). Each lock proves \lean{formal\_claim\_X ↔ paper\_claim\_X}, where \lean{paper\_claim\_X} is transcribed from the corresponding theorem sentence of \lean{papers/fefferman-clay-a-conditional/main.tex} (sha256 \hash{14179f79141f4a2b596f6fc08194be737e2f4bec8235d8f14bfc44175fa939dd}). The English statements of Theorems~\ref{thm:T1}, \ref{thm:T10} and~\ref{thm:T11} above reproduce those three locked sentences, with cross-references renamed. The \lean{truth-gate} workflow run 44047249293 on commit \commit{} reports \lean{TRUTH GATE AUDIT PASS (Fefferman locks)} for these three theorems. Theorems~\ref{thm:T2}--\ref{thm:T9} are not covered by any Truth Gate lock; for them this paper gives the verbatim Lean statement and a reading of it.

\section{Certificate and reproducibility}
\label{sec:repro}
\begin{description}[leftmargin=*,style=nextline]
\item[Repository and commit] \href{https://github.com/BenFrohman/frohmanian-tether-sdiff}{\lean{BenFrohman/frohmanian-tether-sdiff}} at \commit{} (\lean{main} after pull request \#14). Its Lean sources are identical to \prcommit{}, the merge commit of \href{https://github.com/BenFrohman/frohmanian-tether-sdiff/pull/12}{pull request \#12} (merged 17 November 2026, 01:53 CT): \lean{git diff} between the two commits changes only four files under \lean{.github/workflows/} and adds \lean{scripts/collect\_lake\_certificate.sh} and \lean{scripts/finalize\_lake\_certificate.sh}; no \lean{.lean} file, \lean{lakefile.lean}, \lean{lean-toolchain} or \lean{lake-manifest.json} changes.
\item[CI run (kernel build)] ``Lean certificate'' workflow, run \href{https://github.com/BenFrohman/frohmanian-tether-sdiff/actions/runs/44047249311}{44047249311} (event \lean{push}, conclusion \lean{success}). Steps: install elan; record the lake certificate environment; \lean{lake exe cache get}; \lean{lake build}; \lean{lake build Scratch}; source scan for \lean{sorry}/\lean{admit}/\lean{native\_decide}/\lean{axiom}; strict axiom audit; Scratch axiom audit; \lean{bash scripts/verify.sh}; upload of the lake certificate bundle. All succeeded.
\item[Companion runs on the same commit] ``Kernel certificate'' run \href{https://github.com/BenFrohman/frohmanian-tether-sdiff/actions/runs/44047249332}{44047249332} (success; 43 of 43 modules replayed through the kernel checker) and ``truth-gate'' run \href{https://github.com/BenFrohman/frohmanian-tether-sdiff/actions/runs/44047249293}{44047249293} (success).
\item[Manifest] \lean{lake-manifest.json} at \commit{}: sha256 \hash{7d1ab35f9874f0c261e99ae06fe05446a8cec85dd4c9af79b016b73032bc0955}, git blob \hash{0a1387d26aef7e78d6d1b507ee4e95ffa70aa131}. Runs 44047249311 and 44047249332 each upload this file and its sha256; the run log records the checkout of \commit{} and the cache step cloning each dependency at the manifest's revision.
\item[Toolchain] \lean{leanprover/lean4:v4.30.0-rc1}; \lean{Lean (version 4.30.0-rc1, x86\_64-unknown-linux-gnu,} \lean{commit} \hash{714601baf118066cbf3f282361339c6d06665b2a}\lean{, Release)}; \lean{Lake version 5.0.0-src+714601b (Lean version 4.30.0-rc1)}.
\item[elan] \lean{elan 4.2.4 (227caca13 2026-08-25)}, as recorded by runs 44047249311 and 44047249332.
\item[Mathlib] \lean{v4.30.0-rc1}, revision \hash{0c154d67103f74be3a0f2c509f72ccbf5be9f2a7}.
\item[Audits] \lean{SCRATCH AUDIT SUMMARY: theorems audited: 235; clean: 235; violations: 0}; \lean{AUDIT SUMMARY: theorems audited: 1889 …; violations: 0}; source scan: 0 hits; \lean{VERIFY OK}.
\end{description}

\noindent\textbf{\#print axioms}, verbatim from the \lean{lake build Scratch} step of run 44047249311:
\VerbatimInput[fontsize=\scriptsize]{lean/axioms-main.txt}

\noindent\textbf{One rebuild command} (fresh clone; installs elan if absent; uses the committed manifest and the Mathlib cache):
\begin{Verbatim}[fontsize=\scriptsize,breaklines=false]
git clone https://github.com/BenFrohman/frohmanian-tether-sdiff.git && \
  cd frohmanian-tether-sdiff && \
  git checkout 1ea649e495b5e06fafec75a05bb8ec07f13dae80 && \
  bash scripts/verify.sh && lake build Scratch && \
  lake env lean scripts/PrintScratchAxioms.lean
\end{Verbatim}

The certificate bundle deposited with this paper contains the pinned \lean{lake-manifest.json}, \lean{lean-toolchain} and \lean{lakefile.lean} from \commit{}; the GitHub run-log archives and artifact archives of runs 44047249311, 44047249332 and 44047249293 exactly as downloaded from the GitHub API (artifact archive sha256 equal to GitHub's recorded digests); the run, job and artifact metadata; and \lean{SHA256SUMS}.

\section*{License and citation}
Copyright (c) 2026 Benjamin Stanley Frohman. This paper is licensed under CC BY 4.0. The Lean code quoted here and in the repository is licensed under Apache-2.0. Please cite this work and credit Benjamin Stanley Frohman (ORCID 0009-0006-7068-3718) when using it; see the accompanying \lean{NOTICE}.

\section*{Related}
B.\,S.\ Frohman, \emph{Three Pillars of 3D Incompressible Navier--Stokes Theory: Kato Local Existence, Fujita--Kato Small-Data Globality, and the Beale--Kato--Majda Continuation Criterion (While the Roof Remains Open)}, Zenodo, \href{https://doi.org/10.5281/zenodo.24301607}{doi:10.5281/zenodo.24301607}.

\end{document}
